\chapter[Photo and Film: Copy or Creation?]{Do Photography and Film Copy Reality or Create It?}
\section{The Graduation Photograph in a Drawer}
Pick up something that may lie in a corner of a drawer at home: your primary-school graduation photograph.

Its laminated surface is slightly sticky and its corners curl. More than forty people stand in four height-ordered rows, some crouching, some standing, some on benches. You find yourself immediately: third row, fifth from the left, smiling stiffly because sunlight caught your eyes as the photographer counted one, two, three.

The photograph seems utterly uncontroversial. It offers no opinion or argument, merely proof that these people stood on this playground that day in those clothes. If someone asks twenty years later whether everyone attended, you can slap it on the table: see for yourself.

But look longer. You remember morning rain and wet patches, absent from the sunny picture because everyone waited. Your teacher moved your best friend to another row because together you kept laughing. The photographer told everyone to smile, so forty faces adopted one expression simultaneously. A classmate was absent, yet the picture looks complete, and nobody notices the missing person, including you ten years later.

Which is that day: rain, an absence, rearranged people, and commanded smiles, or sunshine, order, and uniform happiness?

The photograph has not lied: what it captured really existed before the lens. But neither has it delivered the day. It has delivered forty faces and a carefully chosen moment of the best weather.

This chapter asks whether photography and film \textbf{copy reality} or \textbf{create reality}. A machine neither imagines nor takes a position or understands rhetoric. Yet behind every image stand people choosing the frame, requesting smiles, and deciding which shot survives. How far do their choices change the reality you see?

\section{Inside a Studio a Century Ago}
Enter a scene with me.

It is around 1900, in a photographic studio in a Chinese treaty port: Shanghai, Tianjin, or Guangzhou.

Lift the curtain and light greets you first. Old studios often occupied top floors, with glass roofs and a glass wall under layers of white curtains. The photographer adjusted them like an invisible orchestra. A high-backed chair stands centrally, with an iron support behind it and a small clamp for the back of your head. Not torture equipment, but help staying still: slow photographic materials required seconds or longer, and a blink or turn could blur you into a white ghost.

An elder seats you and adjusts the support. The photographer disappears under the camera's black cloth and extends a hand: look here, do not move. You stare, hardly daring to shift, as he counts seconds. Before magnesium flash was widespread, photography relied on daylight; sunny days brought business, cloudy days left photographers reading newspapers outside idle studios.

Nearby talk makes you more nervous. Photography arrived in China only decades earlier, and people say the black box steals a person's spirit, fastening a shadow to glass and diminishing the soul. Some pray first, some refuse, some enshrine the resulting picture like another self. Before laughing at superstition, remember the fear. When something can carry away, preserve, copy, and send your likeness to strangers you will never meet, your relationship with your image changes permanently. Today's fears of faces used to train models or inserted into others' videos share the structure of that stolen-soul anxiety.

Photography was a solemn ritual. Families wore their best clothes and sat gravely, not because they could not smile but because these might be their only one or two pictures. A photograph would hang for decades and keep them present after death. What you now shoot and delete in three seconds at a tourist site was entrusted with family memory.

From its arrival in China, photography was therefore never a casual transcription. Roof angles, braces, chosen clothing, managed expressions, and waiting for weather made each photograph a rehearsed ceremony. Why, over a century after the machine appeared, do we still tend to believe a photograph shows exactly how things were?

\section{Can a Machine Lie?}
Set out the conflict before drawing conclusions.

One intuition feels almost like common sense: photographs are machine-made and machines have no motives. People lie because they want you to believe something; cameras only record incoming light. Light reflected from a person, tree, or street passes through lenses and leaves traces on film, now a sensor. This physical chain passes through no human tongue. Courts admit photographs, news uses them as proof, and insurers request accident pictures because of that chain. A mouth can lie; how can a beam of light?

The opposite intuition follows from your graduation picture. The machine does not lie, but says almost nothing. Its operator decides what to photograph or exclude, from what distance and at what second, which shots to retain or delete, crop or brighten. Every decision is human and alters what you see. A camera is a pen: the pen does not lie, but its words depend on the hand holding it.

The real question emerges:

\textbf{Where does an image's evidential force come from?} Honest light, or our trust in a motiveless machine? If light, every photograph should remain reliable wherever light traveled. If trust, how do we recognize its exploitation?

This is no idle discussion. It affects whether a surveillance still can help convict, whether a video can ruin a reputation, which photograph a history book prints, and how to recalibrate your eyes when software can produce convincing events that never happened.

Before continuing, take a position. Someone says, ``There is photographic proof; how could it be false?'' Do you nod, or ask something in return?

\section[Courts and Photographed People]{Courts, Darkrooms, and Photographed People}
At least three groups answer whether photographs are trustworthy. Hear each, first giving the easiest target for ridicule its strongest case.

\textbf{Voice one: photographs can testify.}

State the trusting side fully. It deserves serious attention, because the rest of the chapter may seem stacked against it.

Its reasons have three layers. First, physical cost: light from a scene acts on silver-salt grains in film, making a causal connection to reality. Painting invents freely; photography at least requires something before the lens. Traditional darkroom forgery able to withstand experts was costly and difficult, making photographs harder to fake than speech. Second, oversight of power: images secured many historical truths. Photographs and films at the postwar Nuremberg trials made claims of ignorance untenable. Without images, atrocities might end as conflicting accounts. Photographs are among the powerless person's few weapons; perpetrators can burn documents but not widely dispersed negatives. Third, institutional scrutiny: courts allow challenges and expert examination, and reputable newsrooms verify sources. The trusting side believes not one photograph but this apparatus.

That is a strong case. Reject it only after considering who loses most if photographs cease to count as evidence: those with no other means of making themselves heard.

\textbf{Voice two: manipulation existed from the beginning.}

Skeptics hold thicker case files than most people imagine.

An 1860s standard portrait showed Abraham Lincoln standing grandly beside a table. It was a composite: his head on politician John Calhoun's body, chosen for its dignified pose. After Soviet security chief Yezhov fell and was executed in the 1930s, officials erased him from a photograph beside Stalin at a canal. In 1994, Time darkened football star O. J. Simpson's police photograph for its cover; Newsweek's original alongside it revealed guilt suggested by tonal adjustment. In 1982, National Geographic digitally moved two Giza pyramids closer for a cover, prompting the first public controversy over digital alteration in a major magazine. In 2007, Shaanxi farmer Zhou Zhenglong's purported wild South China tiger photographs excited China before proving to show a tiger poster.

Notice the range: political propagandists, leading magazine designers, and an ordinary farmer. Image manipulation belongs to no single regime or system. It grows from trust in the machine: wherever people believe light cannot lie, someone sets a trap for it.

\textbf{Voice three: the people before the lens often have no voice.}

This asks not whether a picture is true, but who controls its use.

In 1936, Dorothea Lange photographed a mother with children in a California pea-pickers' camp, brow furrowed and hand at her chin. Migrant Mother became a defining Depression image, endlessly reproduced. Lange did not record her name, and the symbol of suffering had no say for decades in how her image circulated. When journalists later found her, she was an impoverished older woman with complex resentment: everyone knew her face, nobody her name. Nanook of the North, the 1922 documentary, similarly introduced an Inuit man to the world under a name that was not his own, performing hunting methods reconstructed according to the director's imagined past.

Each voice holds part of the truth. Moving forward requires not another choice of sides but a tool dividing the broad category of photograph into gates we can inspect separately.

\section{Thinking Bridge: Five Gates of an Image}
Every photograph or video passes through five gates between reality and your eyes, each with a chooser behind it. When an image feels particularly persuasive, open them one by one.

\textbf{Gate one: framing---what enters, what is cut out?}

The frame is a sharp pair of scissors. Photograph a rally's densest corner close up and it becomes a sea of people; a wide elevated view at the same moment may show sparse attendance. Neither fabricates faces, but apparent scale can differ tenfold. Ask what lies three meters outside the frame.

\textbf{Gate two: focal length---how is space compressed?}

The technical term concerns lenses flattening or expanding space. A close wide-angle view exaggerates distance through a crowd, making a dozen people seem to stretch endlessly; a distant telephoto compresses bridge traffic into apparently bumper-to-bumper congestion on a clear road. Faces too acquire different nose sizes and ear spacing. The lens arranges spatial experience. Is this distance your eye's or the lens's?

\textbf{Gate three: exposure and tone---brighter or darker?}

Brighten a face and it seems open; darken it and it seems sinister. Time's Simpson cover worked here. Technically exposure and tone, emotionally mood: what prejudgment does this lighting encourage?

\textbf{Gate four: editing---which slice of time?}

Photography already edits: a continuous day becomes a few hundredths of a second. Keep one of thirty frames and twenty-nine vanish. Retain an athlete's furious expression and the calm half-seconds surrounding it disappear. Film joins moments from different times and places into a seemingly continuous river. Short videos widen this gate: two hours of conversation become the most explosive minute. What happened immediately before and after?

\textbf{Gate five: music and text---the offscreen conductor.}

An image never arrives alone. A crying child captioned lost their home provokes anger; separation anxiety, sympathy; cheerful music can even turn it into comedy. Captions and scores are unseen narrators, changing no pixel yet rewriting meaning. Would you reach the same conclusion with sound muted and title hidden?

These are not five offenses. Each is legitimate technique: without framing no composition, without editing no cinema. The tool does not license calling everything fake. It replaces the unproductive question whether a photograph is true with answerable questions about its choices and where each steers you.

\section[Lu Xun on Photography]{A Century-Old Testimony: Lu Xun on Photography}
Read a short original source. Photography's performance was recognized long before digital images, by one of China's least forgiving observers.

In 1924, Lu Xun wrote On Photography and Suchlike about contemporary studios. A famous passage describes photographing one customer in two costumes and poses, then combining them so two selves appear, one seated and one standing, or one pouring tea for the other. These were double-self pictures. He said customers liked seeing one self seated arrogantly and the other kneeling in abject misery. A composite trick revealed the mind's addiction to hierarchy.

Moving to stage performance, he made an even harsher remark:

\begin{quote}
Our greatest, most enduring, and most widespread Chinese art is men playing women.
\end{quote}
This is satire aimed at an aesthetic fashion he disliked, not necessarily defensible in every respect today. Reducing a performance tradition to a symbol of national character is itself debatable. We quote him not as infallible authority but because he grasped something larger in 1924: \textbf{people often enter studios not to record who they are, but to perform who they wish to be}.

He also mentioned resistance to photography: some believed it stole the spirit and was especially inadvisable during good fortune. Put the two together. Some fear photographs fixing and carrying away the real self; others love them for making a false, grander, higher-status or doubled self more convincing than reality. Fear and fascination point to the same fact: photographs were never mere copies, but material people could manipulate.

A century later, studios have become front-facing phone cameras and double-self pictures filters and beautification. Lu Xun might need to change only the nouns, not rewrite the essay.

\section{One Photograph, Three Explanations}
Return to why photographs count as evidence. Distinguish events---people trust them and courts accept them---from explanations of why, participants' self-understanding such as stolen souls, and judgments of whether trust is rational or naive. Here compare the explanatory level: where does evidential force originate?

\textbf{Explanation one: physics---a trace left by reality.}

French film critic Andre Bazin argued in the mid-twentieth century that photography's power lies in its automatic character: an object impresses itself on film like a fingerprint on glass. Roland Barthes later likewise emphasized proof that something stood before the camera. Unlike drawing or testimony made by human hands, a photograph is made by light's hand. This clear explanation shows why photographs count as evidence while sketches do not. But the chain proves only a scene before the lens, not that it was unstaged or uncomposited. In 1917, two English girls photographed paper fairies in a garden, convincing adults including Arthur Conan Doyle. Light really struck the paper; the causal chain remained intact, but the fairies were false.

\textbf{Explanation two: training---we are taught to believe.}

Perhaps photographs are not naturally credible: societies teach conventions for reading pictures. Late-Qing viewers saw dangerous magic rather than neutrality; modern viewers find old pictures awkward, staged, or fake. The pictures have not changed, only the eyes reading them. Pictures-as-proof would then be learned culture, not instinct. This explains differences across time and cultures, but if training were everything, societies might learn the opposite, treating photographs as inherently suspicious. Instead, basic trust spread rapidly almost everywhere photography arrived. Habit alone does not explain that consistency; photographs seem to possess some distinctive persuasive force.

\textbf{Explanation three: institutions---trust the organization behind it.}

Consider provenance. Identical images on an anonymous account and in a court file carry different weight, not because of pixels but examination, authentication, and cross-questioning. Trust in news photography often means trust in verification and reputation costs; trust in surveillance means confidence in untampered custody. This explains platform-era collapse: millions of institution-detached pictures race through algorithms faster than verification can follow, expiring seeing-is-believing. Yet institutions manipulate too: Time and National Geographic are leading examples. Outsource all trust to them and their collapse may take your capacity for doubt with it.

Each explanation captures something: light was there, conventions trained us, and institutions checked. This is not fence-sitting but a methodological reminder: an explanation claiming everything often discarded something beforehand.

\section[A Third Thing Between Two Shots]{Deeper Challenge: A Third Thing Between Two Shots}
The chapter's weightiest theoretical device comes from an early cinema experiment, a key to all visual narrative.

In the 1920s, Soviet filmmaker Lev Kuleshov conducted an experiment later described by colleague Vsevolod Pudovkin. He joined \textbf{the same neutral close-up} of actor Ivan Mozzhukhin to three images: steaming soup, a child in a coffin, and a woman reclining on a sofa. Viewers praised his acting, reading hunger, grief, and desire in his face, although it never changed by a pixel.

The \textbf{Kuleshov effect} makes a simple, unsettling point: meaning arises not in a single shot but at the \textbf{join} between shots, half supplied by the viewer. Filmmakers recognized a tool, not a defect. Eisenstein developed this into montage theory: two joined shots create something beyond addition. He compared Chinese characters: mouth, \mbox{口}, plus dog, \mbox{犬}, forms bark, \mbox{吠}; water, \mbox{水}, plus eye, \mbox{目}, forms tears, \mbox{泪}. Meaning's unit is the collision between shots, not the shot alone.

In 1925, Eisenstein made Battleship Potemkin to commemorate the 1905 revolution. Its Odessa Steps climax shows firing soldiers, fleeing crowds, a shot mother releasing her pram, and the pram tumbling down endless steps. These minutes became among cinema's most imitated scenes. Historians generally recognize unrest and repression in Odessa in 1905 but consider the steps massacre largely invented. Generations nevertheless remember it as history. Images manufacture memory more thoroughly than words by bypassing someone-told-me and impersonating I-saw-it-myself.

More concealed is \textbf{the viewer's position}. You stand where the camera stands. Low angles make people imposing. Leni Riefenstahl's Triumph of the Will, filmed at the 1934 Nazi rally, used low views and monumental framing to make Hitler a descending giant. Studying this acknowledged propaganda reveals how camera position and editing alone steer emotion without slogans. High angles shrink figures; following someone places you inside their predicament, a killer in horror or detective in mystery. Camera position allocates sympathy. \textbf{A viewpoint is an invisible position taken}.

Documentaries belong in this framework. Their difference from fiction is not editing versus no editing, but claiming a basis in reality versus not making that claim. Documentaries frame, wait, assemble, score, and choose heroes too. Nanook of the North's celebrated walrus hunt used traditional harpoons at the director's request despite hunters already using guns; an igloo was cut open for filming light. These facts do not discredit all documentaries but define them accurately: \textbf{documentary is not unselected reality, but narrative grounded in reality}. Judge whether choices are faithful to its claimed subject, not whether choices exist; without them there is no film.

One easy mistake remains. Learning montage can lead to everything-is-manipulation nihilism. Slow down. Kuleshov's face, soup, and coffin were real; editing created a \textbf{relationship}, not objects from nothing. New meaning at joins does not make every shot false. Equating selection with fabrication destroys the more important ability to distinguish honest narrative from propaganda. Nihilism abandons distinctions rather than deepening them.

\section{Today: The World in Fifteen Seconds}
This photography-old question now races inside your pocket.

First, quantity. Film offered thirty-six exposures, days of processing, and a cost per shot, encouraging thought before pressing. Digital cost fell to zero. Your pocket device can produce more images in a day than a century-old studio in a year. Photography became breathing rather than ritual, and breathing is unthinking. An unexpected consequence appeared in psychologist Linda Henkel's museum experiment: visitors who photographed objects remembered fewer details than those who only looked. The photo-taking-impairment effect suggests that outsourcing memory to the camera makes the brain neglect its own storage. Sometimes photographs \textbf{replace} rather than preserve memory.

Images also occupy memory in reverse. The Odessa pram is Eisenstein's shot, not history, yet now inhabits your mental 1905. Many remembered historical scenes come from film and television; repeated archival footage can overwrite participants' diverse recollections. In the textual era, altering collective memory required burning books; in the visual era, repeating one image can suffice.

Then attention. Film researchers have found Hollywood's average shot length shrinking over decades, from around ten seconds in classic cinema to a few in commercial films. Short videos push further: a complete stimulus in fifteen seconds, music supplying ready-made emotions, editing removing pauses, algorithms recording the extra half-second you linger and serving more. Kuleshov found meaning at joins; platforms find \textbf{watch time}. Montage becomes engineering, aiming less at thought than preventing your thumb from moving. Short videos are not inherently the problem; a few seconds can tell something well. But when everything adopts that rhythm, slow articles, conversations, and thinking time become intolerable.

Finally, the newest gate. Lincoln's composite took a darkroom expert days; a phone now inserts a face into video or fabricates a speech. Forgery costs have fallen from state-propaganda-department scale to a ten-minute school break. Seeing-is-believing has expired, but this does not mean trusting nothing. Add two questions to the five gates: \textbf{source}, where did this first appear and who brought it to me, and \textbf{motive}, what feeling and action does it seek? Finding an original source often takes minutes. Many are deceived not for lack of tools but because the picture is too satisfying to question. Manipulators depend most on willingness to believe, not technology.

\section{For You: When Eyes Can Be Deceived}
Return to the graduation photograph.

You now notice the absent classmate beyond the frame, commanded smiles, and rain waited out. Yet you probably will not tear it up. It remains among that afternoon's few real traces: twelve-year-olds did stand there, and sunlight did strike your eyes.

Across a century, we met soul-fearing customers, composite Lincoln, erased Yezhov, paper fairies, a tumbling pram, and short videos turning three times in fifteen seconds. Two extremes answer one question. One says staging, editing, composites, and fabrication mean trusting nothing. The other warns that suffering and injustice visible only through images would lose their last witnesses, a world without photographs perpetrators would welcome.

So the question for you is:

\textbf{What should believing look like when images can be freely manufactured?}

Raise the threshold until you trust only places your feet reach? Or learn the five gates, source, and motive, giving qualified trust ready for revision? The first feels safe but isolates you from distant worlds. The second takes effort and always risks betrayal. Choose and explain, preferably through an image that personally persuaded or deceived you.

There is no standard answer. But from now on, when your phone's frame lights up, you stand not only behind the lens but among a century of makers and viewers. Machines neither lie nor take responsibility for truth. That responsibility has always been yours.
